% !TEX root = ../master.tex
\section{Kritische Reflexion und Ausblick}
\label{sec:zusammenfassung}

Dieses Kapitel fasst die Ergebnisse der Studienarbeit zusammen und ordnet die Portierung der Python-Lehrbeispiele in TypeScript hinsichtlich Zielerreichung, Nutzen und Grenzen ein. Abschließend werden konkrete Erweiterungs- und Forschungsperspektiven für eine weitergehende, stärker automatisierte Übersetzung diskutiert.

\subsection{Erreichte Ziele und zentrale Ergebnisse}
Ziel des Projekts war es, die in der Lehrveranstaltung \emph{Formale Sprachen} verwendeten, ausführbaren Python-Beispiele in eine TypeScript-basierte Umgebung zu überführen, ohne die fachliche Semantik der Beispiele zu verändern. Dabei standen zwei Aspekte im Vordergrund:

\begin{enumerate}
  \item \textbf{Strukturelle Semantiktreue:} Die übersetzten Implementierungen sollen dieselben Konzepte und Algorithmen ausdrücken (z.\,B.\ Scanner-/Parser-Pipelines, AST-nahe Datenformen, Automatenkonstruktionen).
  \item \textbf{Explizite Typisierung:} Implizite Python-Datenformen sollen in TypeScript als \emph{statisch überprüfbare} Typen modelliert werden, sodass Inkonsistenzen frühzeitig erkannt werden und Schnittstellen präziser dokumentiert sind. \autocite[S.44ff]{vanderkam2019effective}
\end{enumerate}

Als Ergebnis liegt eine TypeScript-Codebasis vor, die die wesentlichen Lehrinhalte aus den Python-Notebooks in einer Umgebung mit statischer Typprüfung reproduzierbar bereitstellt. Der entscheidende Unterschied zur Python-Referenz besteht darin, dass viele zuvor nur informell angenommene Invarianten (z.\,B.\ ``Parser liefert immer (Wert, Restliste)'', ``Transitionstabelle hat wohldefinierte Schlüssel'') nun im Typmodell explizit abgebildet sind und damit bereits durch den Compiler überprüft werden.

\subsection{Mehrwert der statischen Typisierung im Lehr- und Projektkontext}
Die Portierung zeigt, dass TypeScript im Kontext compiler-naher Strukturen (Token, Parserzustände, ParseTrees/ASTs, Automaten) insbesondere dann Vorteile bietet, wenn Datenformen \emph{rekursiv} sind oder wenn viele Teilkomponenten über klar definierte Schnittstellen interagieren. Der Mehrwert lässt sich in drei Punkten bündeln:

\subsubsection{Frühzeitige Fehlererkennung und bessere Schnittstellen}
Durch die Typprüfung werden Fehler, die in Python typischerweise erst bei der Ausführung sichtbar werden, bereits beim Übersetzen erkannt. Beispiele sind vertauschte Rückgabekomponenten, vergessene Fälle in Union-Typen oder inkonsistente Schlüsseltypen in Abbildungen. In diesem Sinne übernimmt das Typsystem eine Rolle als \emph{Prüfinstanz}, die zulässige Zustände eingrenzt und damit das Risiko ``schiefgehender'' Programmpfade reduziert. \autocite[S.1ff]{pierce2002tapl}

\subsubsection{Dokumentation durch Typen und Wartbarkeit}
Typdefinitionen (insbesondere bei Union-/Tupel-Modellen) wirken als präzise Dokumentation. Für Lehrbeispiele ist dies didaktisch relevant: Studierende sehen nicht nur den Algorithmus, sondern auch dessen erwartete Datenformen explizit. Zudem werden spätere Erweiterungen (z.\,B.\ neue Tokenarten, zusätzliche AST-Formen) systematisch durch den Compiler begleitet, weil an betroffenen Stellen Typfehler auftreten, bis alle notwendigen Anpassungen erfolgt sind. \autocite[S.44ff]{vanderkam2019effective}

\subsubsection{Strukturelle Typisierung als Brücke zu Python-Datenmodellen}
Da TypeScript strukturell typisiert ist, können viele Python-nahe Datenmodelle direkt als Records, Tupel und Unions modelliert werden, ohne eine nominelle Klassenhierarchie einzuführen. Das ist für die Portierung besonders hilfreich, weil die Python-Notebooks häufig \emph{datenorientiert} (statt objektorientiert) formuliert sind. \autocite[S.44ff]{vanderkam2019effective}

\subsection{Einordnung der Transformation als compiler-nahe Übersetzung}
Die in Kapitel~4 beschriebene Transformation lässt sich konzeptionell als vereinfachte, syntaxgerichtete Übersetzung verstehen, bei der aus einer Quellrepräsentation schrittweise eine Zielrepräsentation erzeugt wird. In compiler-theoretischer Terminologie entspricht dies einer Kombination aus:
(i) syntaktischer Analyse (Parser), (ii) semantischer Modellierung/Transformation (AST-Transformation) und (iii) Codegenerierung. \autocite[S.304ff]{aho2006compilers}

Die praktische Portierung ist dabei weniger als \emph{vollautomatische Übersetzung} zu verstehen, sondern als \emph{systematische Rekonstruktion} der Programme unter Beibehaltung der algorithmischen Kernaussage. Die wichtigste Zusatzleistung entsteht nicht durch neue Algorithmen, sondern durch die explizite Formalisierung zuvor impliziter Python-Datenformen in Form von TypeScript-Typen.

\subsection{Grenzen der Portierung}
Trotz der erzielten Semantiktreue bestehen Grenzen, die aus den unterschiedlichen Laufzeit- und Sprachmodellen resultieren:

\begin{itemize}
  \item \textbf{Python-spezifische Dynamik:} Python erlaubt sehr flexible Datenformen (z.\,B.\ heterogene Listen, spontane Strukturänderungen, ``duck typing'' zur Laufzeit). Eine 1:1-Übernahme dieser Flexibilität ist in TypeScript entweder nicht sinnvoll oder würde zur Verwässerung des Typmodells (z.\,B.\ durch \texttt{any}) führen. \autocite[S.44ff]{vanderkam2019effective}
  \item \textbf{Unterschiede im Zahlenmodell:} Python unterscheidet \texttt{int} und \texttt{float}, TypeScript fasst dies in \texttt{number} zusammen. Semantisch problematisch wird dies bei sehr großen Ganzzahlen oder exakter Ganzzahlarithmetik.
  \item \textbf{Gleichheit/Schlüsselstrukturen:} Python kann komplexe, unveränderliche Strukturen (z.\,B.\ \texttt{frozenset}, \texttt{tuple}) als Schlüssel in \texttt{dict} verwenden. In TypeScript muss diese Semantik durch geeignete Container-/Schlüsselrepräsentationen explizit nachgebildet werden, da \texttt{Map}/\texttt{Set} referenzbasiert arbeiten.
\end{itemize}

Diese Grenzen sind nicht als Mängel der Portierung zu verstehen, sondern als Konsequenz der unterschiedlichen Designziele beider Sprachökosysteme. In der Arbeit wurde daher konsequent priorisiert, \emph{Typpräzision und Semantik} zu erhalten, statt Python-Dynamik um jeden Preis zu simulieren.

\subsection{Ausblick: Technische Erweiterungen und Forschungsfragen}
Aus dem Projekt ergeben sich mehrere naheliegende Erweiterungen, die sowohl didaktisch als auch technisch relevant sind:

\subsubsection{(Teil-)Automatisierte Typableitung und Übersetzung}
Ein nächster Schritt wäre die stärkere Automatisierung der Übersetzung: Anstatt Typabbildungen manuell zu rekonstruieren, könnten Regeln für die Ableitung von TypeScript-Typen aus Python-Annotationen und aus beobachteten Datenformen systematisiert werden. In compiler-naher Perspektive entspricht dies einer Erweiterung der ``semantischen Regeln'' der Übersetzung. \autocite[S.304ff]{aho2006compilers}

\subsubsection{Systematische Korrektheitsabsicherung}
Kapitel~5 validiert die Semantik primär empirisch (Tests, Vergleichsausführungen). Als Ausblick bietet sich an, Teilaspekte formaler zu fassen, etwa für ausgewählte Kernfunktionen:
\[
\forall P \in \mathcal{P}:\ \llbracket P \rrbracket_{\text{Python}} = \llbracket f(P) \rrbracket_{\text{TS}}
\]
In der Praxis könnte dies durch stärkeres Property-based-Testing oder durch formale Spezifikationen einzelner Komponenten angenähert werden. \autocite[S.15ff]{aho2006compilers}

\subsubsection{Didaktische Aufbereitung als Toolchain}
Langfristig kann die TypeScript-Version als eigenständige Toolchain für Lehrzwecke ausgebaut werden (z.\,B.\ als NPM-Paket oder Web-Demo), sodass Studierende direkt in einer einheitlichen Umgebung experimentieren können. Besonders attraktiv ist dabei die Möglichkeit, Typfehler als unmittelbares Feedback in IDEs zu nutzen (Autocomplete, Refactoring, statische Hinweise).

\subsubsection{Erweiterung des Sprachumfangs}
Die übersetzten Beispiele decken einen klaren Ausschnitt ab (formale Sprachen, Parsing, Automaten, AST-nahe Ausdrücke). Eine Erweiterung auf weitere Python-Konstrukte (z.\,B.\ komplexere Ausdrucksformen, zusätzliche Datentypen, mehr Standardbibliothek) müsste bewusst abwägen, welche Dynamik sinnvoll typisierbar ist und wo die Domäne eher eine abstrahierte Modellierung verlangt.

\subsection{Schlussbemerkung}
Die Arbeit zeigt, dass die Portierung von Python-Lehrbeispielen in TypeScript nicht nur eine syntaktische Übertragung ist, sondern eine \emph{Modellierungsaufgabe}: Implizite Datenformen werden explizit, informelle Invarianten werden zu Typen, und die Compilerprüfung wird zu einem integralen Teil der Qualitätssicherung. Damit entsteht eine robuste Grundlage, auf der die behandelten Inhalte aus dem Bereich formaler Sprachen konsistent, wartbar und didaktisch nutzbar bereitgestellt werden können. \autocite[S.44ff]{vanderkam2019effective}